% Part: first-order-logic
% Chapter: syntax-and-semantics
% Section: subformulas

\documentclass[../../../include/open-logic-section]{subfiles}

\begin{document}

\olfileid{fol}{syn}{sbf}

\olsection{\printtoken{P}{subformula}}

\begin{explain}
اغلب سودمند است دربارهٔ !!{formula}sی سخن بگوییم که !!a{formula} داده‌شده
را «می‌سازند». آن‌ها را \emph{!!{subformula}s} آن فرمول می‌نامیم. هر
!!{formula}، !!a{subformula} خودش نیز به شمار می‌آید؛ زیر‌فرمولی از $!A$
که خود $!A$ نباشد، \emph{!!{subformula} سره} نامیده می‌شود.
\end{explain}

\begin{defn}[!!^{subformula} بلافصل]
اگر $!A$ یک !!a{formula} باشد، \emph{!!{subformula}sی بلافصل} $!A$
به‌صورت استقرایی چنین تعریف می‌شوند:
\begin{enumerate}
\item !!{formula}sی اتمی هیچ !!{subformula} بلافصلی ندارند.

\tagitem{prvNot}{\indcase{!A}{\lnot !B}{تنها !!{subformula} بلافصل
    $\indfrm$، فرمول~$!B$ است.}}{}

\item \indcase{!A}{(!B \ast !C)}{!!{subformula}sی بلافصل $\indfrm$،
  فرمول‌های $!B$ و $!C$ هستند ($\ast$ هر یک از رابط‌های دوجایگاهی است).}

\tagitem{prvAll}{\indcase{!A}{\lforall[x][!B]}{تنها !!{subformula} بلافصل
    $\indfrm$، فرمول~$!B$ است.}}{}

\tagitem{prvEx}{\indcase{!A}{\lexists[x][!B]}{تنها !!{subformula} بلافصل
    $\indfrm$، فرمول~$!B$ است.}}{}
\end{enumerate}
\end{defn}

\begin{defn}[!!^{subformula} سره]
اگر $!A$ یک !!a{formula} باشد، \emph{!!{subformula}sی سرهٔ} $!A$
به‌صورت بازگشتی چنین تعریف می‌شوند:
\begin{enumerate}
\item !!{formula}sی اتمی هیچ !!{subformula} سره‌ای ندارند.

\tagitem{prvNot}{\indcase{!A}{\lnot !B}{!!{subformula}sی سرهٔ
    $\indfrm$ عبارت‌اند از~$!B$ همراه با همهٔ !!{subformula}sی سرهٔ~$!B$.}}{}

\item \indcase{!A}{(!B \ast !C)}{!!{subformula}sی سرهٔ $\indfrm$
  عبارت‌اند از $!B$ و $!C$، همراه با همهٔ !!{subformula}sی سرهٔ $!B$
  و همهٔ زیر‌فرمول‌های سرهٔ~$!C$.}

\tagitem{prvAll}{\indcase{!A}{\lforall[x][!B]}{!!{subformula}sی سرهٔ
    $\indfrm$ عبارت‌اند از~$!B$ همراه با همهٔ !!{subformula}sی سرهٔ~$!B$.}}{}

\tagitem{prvEx}{\indcase{!A}{\lexists[x][!B]}{!!{subformula}sی سرهٔ
    $\indfrm$ عبارت‌اند از~$!B$ همراه با همهٔ !!{subformula}sی سرهٔ~$!B$.}}{}
\end{enumerate}
\end{defn}

\begin{defn}[!!^{subformula}]
!!{subformula}sی $!A$ عبارت‌اند از خود $!A$ همراه با همهٔ
!!{subformula}sی سرهٔ آن.
\end{defn}

\begin{explain}
به تفاوت ظریف میان شیوهٔ تعریف !!{subformula}sی بلافصل و
!!{subformula}sی سره توجه کنید. در حالت نخست، !!{subformula}sی بلافصل
فرمول~$!A$ را برای هر صورت ممکن~$!A$ مستقیماً تعریف کرده‌ایم. این تعریفی
صریح و حالت‌به‌حالت است، و حالت‌ها با تعریف استقرایی مجموعهٔ !!{formula}s
هم‌ساخت‌اند. در حالت دوم نیز از شیوهٔ تعریف مجموعهٔ همهٔ !!{formula}s
پیروی کرده‌ایم، اما در هر حالت، !!{subformula}sی سرهٔ !!{formula}sی
کوچک‌تر $!B$ و $!C$ را نیز، افزون بر خود این !!{formula}s، گنجانده‌ایم.
این امر تعریف را
\emph{بازگشتی} می‌کند. به‌طور کلی، تعریف تابعی روی مجموعه‌ای با تعریف
استقرایی (در مورد ما، !!{formula}s) هنگامی بازگشتی است که حالت‌های تعریف
تابع از خود تابع استفاده کنند. بااین‌همه، برای خوش‌تعریف‌بودن باید مطمئن
شویم که فقط مقادیر تابع را برای آرگومان‌هایی به کار می‌بریم که «پیش از»
آرگومان در حال تعریف می‌آیند---در مورد ما، هنگام تعریف «!!{subformula}
سره» برای $(!B \ast !C)$ فقط از !!{subformula}sی سرهٔ !!{formula}sی
«پیشین» $!B$ و $!C$ استفاده می‌کنیم.
\end{explain}

\begin{prop}
\ollabel{prop:subfrm-trans}
فرض کنید $!B$ زیر‌فرمول $!A$ و $!C$ زیر‌فرمول $!B$ باشد. آنگاه $!C$
زیر‌فرمول $!A$ است. به بیان دیگر، رابطهٔ زیر‌فرمول تعدی‌پذیر است.
\end{prop}

\begin{prob}
\olref[fol][syn][sbf]{prop:subfrm-trans} را اثبات کنید.
\end{prob}

\begin{prop}
\ollabel{prop:count-subfrms}
فرض کنید $!A$ فرمولی با $n$ رابط و سور باشد. آنگاه $!A$ حداکثر
$2n+1$ زیر‌فرمول دارد.
\end{prop}

\begin{prob}
\olref[fol][syn][sbf]{prop:count-subfrms} را اثبات کنید.
\end{prob}

\end{document}
